\chapter{MuJoCo下载安装}
\section{C/C++安装}
下载路径https://github.com/google-deepmind/mujoco/releases
由于我们是初学者，并不需要开发或修改核心MuJoCo代码，所以我们使用预编译的二进制文件即可。

下载并解压缩，可以得到以下目录
\begin{itemize}
	\item[bin:] 动态链接库、可执行文件等
	\item[include:] 使用MuJoCo所需要的头文件
	\item[lib:] 使用MuJoCo所需要的头文件
	\item[model:] 模型文件汇总
	\item[sample:]
	\item[simulate:]  
\end{itemize}
在bin子目录下运行:
\begin{lstlisting}
	./simulate ../model/humanoid/humanoid.xml
\end{lstlisting}
如果能够正常运行，则表示环境没什么问题

我们需要OpenGL渲染，所以安装GLFW库
\begin{lstlisting}
	sudo apt install libglfw3-dev
\end{lstlisting}

\section{Python安装}
为了不污染本地的Python环境，我们安装Anaconda，并通过conda创建虚拟环境。(如果和我一样污染了也无所谓那就算了，跳过环境配置)

环境配置：
\begin{lstlisting}
	conda create --name mujoco-sim python=3.8
\end{lstlisting}

在mujoco-sim环境下：
\begin{lstlisting}
	pip install mujoco
\end{lstlisting}

然后在model文件下运行：
\begin{lstlisting}
	python -m mujoco.viewer --mjcf=./humanoid/humanoid.xml
\end{lstlisting}
如果能够正常运行，则表示环境没什么问题